\usepackage{byrne}
\usepackage{etoolbox}
\usepackage{ccicons,bxtexlogo,listings}
\usepackage{tikz}
\usepackage{comment}
\usepackage{fontspec}
\usepackage{lineno}
\usepackage{ebgaramond}
%\usepackage{microtype} % Improves line breaking and spacing
\usepackage{titling}
\usepackage{fancyhdr}
\usepackage{anyfontsize}
\usepackage{titlesec}
\usepackage[cmyk]{xcolor}
\selectcolormodel{cmyk}
%\usepackage{hyperref}


% 1. Format for Unnumbered Chapters (\chapter*)
\titleformat{name=\chapter,numberless}[block]
  {\normalfont\bfseries\Large}{}{0pt}{} % Large font size (17.28pt in 12pt document)

% 2. Format for Sections
\titleformat*{\section}{\normalfont\bfseries\large}

% 1. Remove the "Chapter Sink" (vertical drop)
\setlength{\beforechapskip}{-3\baselineskip} 
% Use a negative value (e.g., -20pt or -1cm) to pull it even higher if needed.

% 2. Adjust space between "Proposition X" and Enunciation
\setlength{\midchapskip}{10pt} 

% 3. Adjust space between "Proposition X" and the text body
\setlength{\afterchapskip}{20pt}

%%%% set format for 7x10 
\usepackage[
paperwidth=7in, 
paperheight=10in, 
inner=1.25in,
outer=0.75in, 
top=1in, 
bottom=1in
\setlength{\headsep}{15pt}
]{geometry} 

%%%% sets format for 8.5x11
%\usepackage[
%paperwidth=8.5in, 
%paperheight=11in, 
%inner=1in,
%outer=1in,
%top=1in,
%bottom=1in,
%headsep=15pt
%]{geometry}


\pagestyle{fancy}
\setlength{\parskip}{0pt}
\setlength{\headheight}{15pt}
\setlength{\parindent}{0pt}  % Set auto-indent to 0 point

% Define header/footer styles
\fancyhf{}
\fancyhead[LE,RO]{\leftmark}  % chapters (Def, CN, post, prop# in header
\fancyfoot[C]{\thepage} % page numbers centered at bottom

% Format footnotes (removes indent at start of footnote)
\makeatletter
\renewcommand\@makefntext[1]{%
  \noindent\makebox[1em][l]{\@thefnmark}#1}
\makeatother

% Number every nth line
\modulolinenumbers[1]

% Main Text Toggles
\label{Toggles}

% Set EB Garamond as the main font
\setmainfont{EB Garamond}

% Define line numbering toggle
\newtoggle{lineNumbering}

% Define text version toggle Byrnified HEATH original or RAUM-modified Heath
\newtoggle{heath} % Byrnified Heath text
\newtoggle{raum} % Raumified Heath
\newtoggle{heath1} % Byrnified Heath text without paragraph formatting
% Set values of toggles (uncomment one to enable it)
\toggletrue{heath}
\togglefalse{raum}
\togglefalse{heath1}

% Define and set toggle for pictures to be included in definitions or not
\newtoggle{defpics}
\toggletrue{defpics}

% Define text version toggle
\newtoggle{inline}

% Set line numbering on or off
\toggletrue{lineNumbering}

% Set inline drawings toggle (comment it out to turn off inline)
\toggletrue{inline}

\label{MetaPost Definitions}

\def\mpPre{
	u:=1cm;
    lineWidth := 2pt; % default is 2pt - increase for thicker lines in drawings
	textLabels := true;
	textLabelScaleFactor := 1; % set label size for diagrams
	textLabelAvSize := abs(ulcorner(textLabelRender("A")) - llcorner(textLabelRender("A")));
	angleSize := 1cm;
	compensateLineLabels := false; % this one removes vertical compensation for text labels for line segments
	lineLabelsOnTop := false; % this one moves line segment labels to the bottom
        angleModeFallback := true;
	}
 
\def\mpPost{
	textLabels := false;
    lineWidth := 2pt; % default is 2pt - increase for thicker lines in drawings
    textLabelScaleFactor := 1; % set label size for inline lines
	}

\label{Commands, Parameters}

\label{Justification Command}
% Robust command for formatting justifications (e.g. "C.N. 1")
 
\newcommand{\justification}[1]{%
  \unskip                       % 1. Remove space before command
  \hfill                        % 2. Push to the right
  \penalty100                   % 3. "If you must break the line, do it here"
  \hskip 1em plus 1fill         % 4. The "Smart Spring"
  \mbox{[\textit{#1}]}%         % 5. The justification (Italicized, normal size)
  \par                          % 6. Finish the line
}

\newcommand{\justificationNew}[1]{%
  \unskip                % Remove extra space
  \hspace*{2em}          % Minimum gap before the note
  \nobreak\hfill         % Try to stay on this line and push right
  \allowbreak            % If it's too tight, allow a break here
  \hspace*{\fill}        % If it breaks, push the note to the right on the new line
  \mbox{[\textit{#1}]}% The note itself
  \par                   % End the line/para
}

%%%%%
% Command for centering main text (#1), right-flushing the justification (#2)

\usepackage{calc}
\newlength{\textlength}
\newlength{\justlength}

\newcommand{\centeredjust}[2]{%
	%\unskip
    % 1. Measure the text and the justification
    \settowidth{\textlength}{#1}
    \settowidth{\justlength}{[\textit{#2}]}
    
    % 2. Logic: If (text + justification + buffer) > linewidth, drop to next line
    \ifdim\dimexpr\textlength+\justlength+1em\relax > \linewidth
        {\centering #1 \par}
        \hspace*{\fill}[\textit{#2}]\par
    \else
        % 3. If it fits, use the "Balanced Ghost" method to keep text centered
        \noindent\hspace*{\fill}\makebox[0pt][l]{\phantom{[\textit{#2}]}}#1%
        \hspace*{\fill}\llap{[\textit{#2}]}\par
    \fi
}
%%%%%

\def\pointlabel#1{\textsf{\textbf{#1}}} % This one determines the style of text labels both on the main diagram and in the text

\newdimen\textLabelShift % This thing allows you to move labeled line segments up and down to make the labels appear like normal characters

\textLabelShift = -1.6pt % This value works for 12pt EB Garamond, but might not work with a different font/size

\newdimen\circleSpace % Space between inline circles and the letters

\circleSpace = 2pt % This value works for 12pt EB Garamond, but might not work with a different font/size

\newdimen\pointShift % This allows you to move the inline points backward (negative) or forward (positive) to fine tune the spacing -- the \pointlabel function seems to usually insert a space before points. 

\pointShift = -3.5pt % This value works for 12pt EB Garamond, but might not work with a different font/size

\newdimen\lineShift % This allows you to move the inline lines backward (negative) or forward (positive) to fine tune the spacing -- the \pointlabel function seems to usually insert a space before points. 

\lineShift = -5pt % This value works for 12pt EB Garamond, but might not work with a different font/size

\newdimen\inlineLength % this is the length of the inline lines. 

\inlineLength = 0.55cm % Length of inline lines should be ~0.046 * fontsize (it's 0.55cm for 12pt)

\DeclareDocumentCommand{\drawUnitLineAlt}{o m O{0.55cm}}{% This is an alternative macro for inline line segments. Optional input is to rename the line
    \iftoggle{inline}{
        \hspace{\lineShift}
        \advance\pictOffsetBottom by \textLabelShift
	    \advance\pictOffsetTop by \textLabelShift
    	\IfNoValueTF{#1}%
    		{\drawFromCurrentPicture[top][uline#2]{%
    			textLabels := true;%
    			startAutoLabeling;%
    			draw byNamedCompoundLine(#3, 0)(#2);
    			stopAutoLabeling;%
    			textLabels := false;%
    			}}%
    		{\drawFromCurrentPicture[top][uline#2]{%
    			textLabels := true;%
    			startAutoLabeling;%
    			draw byNamedLineSeqByPoint(1)(#2);%
    			stopAutoLabeling;%
    			textLabels := false;%
    			}}%
	    \advance\pictOffsetBottom by -\textLabelShift%
	    \advance\pictOffsetTop by -\textLabelShift}%
        {%
            \IfNoValueTF{#1}%
                {\pointlabel{#2}%
                }%
                {\pointlabel{#1}}%
        }%
}%

\DeclareDocumentCommand{\drawAngleAlt}{o m}{% This alternative in-line angle command draws the named angle, then puts the name (or the optional argument) in bold. 
    \iftoggle{inline}{%
        \drawAngle{#2}%
	\IfNoValueTF{#1}%
	{\pointlabel{#2}}%
	{\pointlabel{#1}}%
	}%
    {%
        \pointlabel{#2}%
    }%
}%

 \DeclareDocumentCommand{\drawPolygonAlt}{o m}{% This alternative in-line polygon command draws the named polygon, then puts the name (or the optional argument) in bold. 
    \iftoggle{inline}{%
        \drawPolygon{#2}%
	\IfNoValueTF{#1}%
    	{%
            \pointlabel{#2}}%
    	{%
            \pointlabel{#1}}%
    	}%
    {%
        \pointlabel{#2}%
    }%
}%

%% Define bycolors for inline circles. Other uses of colors are defined in byrne.mp and byrnebook.cls
% RGB equivalents
\definecolor{byblue}{rgb}{.15, .35, .6}
\definecolor{byred}{rgb}{.85, .3, .1}
\definecolor{byyellow}{rgb}{.95, .7, 0.1}
\definecolor{bygreen}{rgb}{0.122, 0.529, 0.122}
\definecolor{byblack}{rgb}{0, 0, 0}
\definecolor{bygrey}{rgb}{.8, .8, .8}
\definecolor{bywhite}{rgb}{1, 1, 1}

% CMYK equivalents
%\definecolor{byred}{cmyk}{0.0, 0.88, 0.88, 0.15}
%\definecolor{byblue}{cmyk}{0.75, 0.417, 0.0, 0.4}
%\definecolor{byyellow}{cmyk}{0.0, 0.263, 0.895, 0.05}
%\definecolor{byblack}{cmyk}{0, 0, 0, 1}
%\definecolor{bygrey}{cmyk}{0.0, 0.0, 0.0, 0.2}

%%%%% drawCircleAlt
\DeclareDocumentCommand{\drawCircleAlt}{o m}{%
  \iftoggle{inline}{%
    \tikz[baseline=-1ex]{
      \draw[line width=1.6pt, draw=#1] (0,0) circle [radius=0.5em];
    }%
    \hspace{\circleSpace}\pointlabel{#2}
  }{%
    \pointlabel{#2}%
  }%
}

%%%%% ^drawCircleAlt

\def\drawPointAlt#1{%
    \iftoggle{inline}{%
        \drawPointM{#1}\hspace{\pointShift}\pointlabel{#1}%
        }%
    {%
    \pointlabel{#1}%
    }%
}%
%\def\drawAngleAlt#1{\drawAngle{#1}\pointlabel{#1}}

\DeclareDocumentCommand{\drawLineAlt}{o m}{% Alternative inline drawing command to draw connected sequence of lines (e.g. for triangles or polygons in general)
    \iftoggle{inline}{%
        \drawLineByPoints{#2}%
    \IfNoValueTF{#1}%
    {\pointlabel{\directlua{cleanPolygonName("#2")}}}%
    {\pointlabel{#1}}%
    }%
    {%
        \pointlabel{#2}%
    }%
}%

\def\RemoveCommas#1,{%
  \ifx\relax#1\relax% Check if we've reached the end
  \else%
    \textbf{#1}% Print in bold without a comma
    \expandafter\RemoveCommas% Continue with the next item
  \fi%
}%



%%%%%%%%%%%%%%%%%%%%%%%%%%%%%%%%%%%%%%%%

\lstset{
language=MetaPost,
alsolanguage=TeX,
numbers=none,
basicstyle=\ttfamily\scriptsize
}

\directlua{
function cleanPolygonName(s)
	local resultingName = s:gsub ("[, ]","")
	resultingName = resultingName:gsub ("-","minus")
	tex.print(resultingName)
end
}


%%%%%%%%%%%%%%%%%%%%%%%%%%%%%%%%%%%%%%%%%%%%%%%%%%%%%%%%%%%%%%%%%%%%
% The following is the original \drawLineAlt command, which had no optional argument to rename the figure
% \def\drawLineAlt#1{
%     \iftoggle{inline}{
%     \drawLineByPoints{#1}\pointlabel{\directlua{cleanPolygonName("#1")}}
%         }
%     {
%         \pointlabel{#1}
%     }
% }
%%%%%%%%%%%%%%%%%%%%%%%%%%%%%%%%%%%%%%%%%%%%%%%%%%%%%%%%%%%%%%%%%%%%

%%%%%%%%%%%%%%%%%%%%%%%%%%%%%%%%%%%%%%%%%%%%%%%%%%%%%%%%%%%%%%%%%%%%
%\def\drawPolygonAlt#1{\drawPolygon{#1}\pointlabel{#1}}
%drawPolygon[verticalAlignment][PolygonNewName] \{Poly\-gon\-Na\-me\}}
%%%%%%%%%%%%%%%%%%%%%%%%%%%%%%%%%%%%%%%%%%%%%%%%%%%%%%%%%%%%%%%%%%%%

%%%%%%%%%%%%%%%%%%%%%%%%%%%%%%%%%%%
% DBM attempt at new drawPolygonAlt command
\DeclareDocumentCommand{\drawMergedPolygonAlt}{o m}{% This is an alternative macro for polygons, allowing for merged polygons
    \iftoggle{inline}{%
        \drawFromCurrentPicture[middle][polygon#2]{%
            startAutoLabeling;%
            draw byNamedPolygon(#2);%
            stopAutoLabeling;%
            \pointlabel{#1}%
            }%
        {%
        \pointlabel{#1}%
        }%
    }%
}%
%%%%%%%%%%%%%%%%%%%%%%%%%%%%%%%%%%%


%%%%%%%%%%%%%%%%%%%%%%%%%%%%%%%%%%%%%%%%%%%%%%%%%%%%%%%%%%%%%%%%%%%%
% KDP 7x10 Left-Right Spread Formatting for Propositions (Fail-Safe)
%%%%%%%%%%%%%%%%%%%%%%%%%%%%%%%%%%%%%%%%%%%%%%%%%%%%%%%%%%%%%%%%%%%%
\usepackage{afterpage}

% Global state tracking
\newtoggle{inproposition}
\togglefalse{inproposition}
\newcommand{\currentpropdiagram}{}

% Automatically inserts the diagram on subsequent left (even) pages if text overflows
\newcommand{\insertEvenPageDiagram}{%
  \iftoggle{inproposition}{%
    % REMOVED the initial \clearpage that was breaking line-numbering and paragraph flow!
    \begin{center}
      \vspace*{\fill}
      \currentpropdiagram
      \vspace*{\fill}
    \end{center}
    \clearpage
    % Re-queue this check for the next page in case text overflows yet again
    \afterpage{\insertEvenPageDiagram}%
  }{}%
}

\newcommand{\startProposition}[2]{%
  \clearpage
  % Force an even (left-facing) page for the primary diagram
  \ifodd\value{page}\hbox{}\clearpage\fi 
  
  % Save the diagram code and set the active proposition state
  \renewcommand{\currentpropdiagram}{#1}
  \toggletrue{inproposition}
  
  % 1. Render the initial diagram on the left-facing page
  \begin{center}
    \vspace*{\fill}
    #1
    \vspace*{\fill}
  \end{center}
  \clearpage % Move cleanly to the right-facing (odd) page
  
  % 2. Defer queuing the even-page diagram until the actual text paragraph begins.
  \everypar{\everypar{}\afterpage{\insertEvenPageDiagram}}%
  
  % 3. Typeset the text block
  #2
  
  % 4. Clean up state before shipping out the final page to prevent trailing loops
  \togglefalse{inproposition}
  \clearpage
}

% Create Header Style
% 1. Create a global variable for the right-hand header
\newcommand{\currentpropheader}{Euclid in Color}

% 2. Define a native memoir page style (bypasses standard LaTeX marks)
\makepagestyle{euclidstyle}
\makeevenhead{euclidstyle}{\fontsize{12}{14}\selectfont Euclid in Color}{}{}
\makeoddhead{euclidstyle}{}{}{\fontsize{12}{14}\selectfont\rightmark}
\makeevenfoot{euclidstyle}{}{\thepage}{}
\makeoddfoot{euclidstyle}{}{\thepage}{}

% 3. Activate the style globally
\pagestyle{euclidstyle}